\chapter{技术方案设计}
\label{chap:design}

本章重点讨论 Agent 方法层设计：如何在受控边界内使用 Agent 完成检测任务规划，如何将既有原子检测能力工具化，如何把多张图像、多类缺陷、多组区域和人工复核意见逐层整合为项目级评估报告。整体上，本系统采用“工程状态机约束 Agent，Agent 补足语义决策能力”的技术方案，使智能能力服务于确定性业务流程。

\section{关键术语}
\label{sec:design-key-terms}

为了避免后续表述产生歧义，本章首先对使用的关键术语进行说明。如\cref{tab:design-key-terms} 所示这里的术语并非孤立概念，而是对应系统中真实存在的设计对象、服务节点和数据结构。

\begin{table}[!htb]
  \caption{技术方案设计关键术语}\label{tab:design-key-terms}
  \centering
  \begin{tabular}{@{}p{0.22\linewidth}p{0.68\linewidth}@{}} \toprule
    \textbf{术语} & \textbf{含义} \\ \midrule
    Agent 自主任务调度方法 & 由分类 Agent 根据幕墙图像中的材质和缺陷线索，输出受控检测任务代码列表，再由业务系统创建和调度后续检测子任务的方法 \\
    LLM 分层信息整合策略 & 将原子检测结果、图像级总结、人工复核和项目上下文分层组织，逐级生成图像级总结和项目级评估报告的策略 \\
    原子检测能力 & 面向某一类幕墙材质或缺陷的最小可调度检测能力，包括金属锈蚀、石材裂缝、石材污渍、玻璃平整度和玻璃爆裂五类 \\
    任务代码 & 用于连接 Agent 决策与工程调度的有限动作集合，即 \texttt{corrosion}、\texttt{crack}、\texttt{stain}、\texttt{flatness} 和 \texttt{spalling} \\
    受控 Agent & 只能在固定业务节点被调用、只能输出固定结构、不能直接写库或推进业务状态的 Agent \\
    检测主任务 & 以单张项目图像为单位创建的检测业务记录，串联分类、子检测、图像级总结和最终状态 \\
    检测子任务 & 由某个原子检测能力执行的具体任务，保存该类检测结果、区域指标、产物信息和执行状态 \\
    图像级总结 & 针对单张图像的结构化检测结果生成的专业文字总结，用于前端展示和项目报告上游输入 \\
    项目级评估报告 & 综合项目基础信息、图像组、图像检测结果、图像总结和人工复核意见形成的建筑幕墙韧性评估结论 \\
    人机协同复核 & 用户对 Agent 和模型结果进行准确性判断与补充评论，并将复核信息纳入项目级报告上下文的机制 \\ \bottomrule
  \end{tabular}
\end{table}

\section{总体方法框架}
\label{sec:design-method-overview}

本文的总体方法框架面向\cref{chap:requirements}提出的系统需求，重点回答两个问题：第一，如何把项目级图像资产、多材质幕墙图像和用户真实评估目标组织为可执行的检测任务；第二，如何把原子检测事实、图像语义、人工复核和项目上下文整合为建筑级韧性评估报告。与单纯的检测模型调用不同，本文方法不把智能能力集中在某一个端到端模型中，而是将系统拆分为“工程约束、数据证据、原子检测、Agent 调度、LLM 整合、质量治理和目标输出”等多个层次，使每一层都有明确职责。

\begin{figure}[!htb]
  \centering
  \resizebox{\linewidth}{!}{%
    \begin{tikzpicture}[
      font=\scriptsize,
      >=stealth,
      region/.style={draw, rounded corners=4pt, line width=0.65pt, align=center},
      layerTitle/.style={font=\bfseries\scriptsize, align=left, anchor=west},
      sideTitle/.style={font=\bfseries\scriptsize, align=center},
      cell/.style={draw, rounded corners=2pt, align=center, minimum height=0.40cm, line width=0.42pt, inner sep=1.7pt},
      note/.style={font=\tiny, align=center},
      legend/.style={font=\tiny, align=center},
      arrow/.style={->, line width=0.58pt},
      dashedarrow/.style={->, dashed, line width=0.55pt},
      roofbox/.style={draw=purple!65!black, fill=purple!7, line width=0.70pt},
      demandbox/.style={region, draw=blue!55!black, fill=blue!4},
      outputbox/.style={region, draw=magenta!55!black, fill=magenta!4},
      llmbox/.style={region, draw=green!50!black, fill=green!5},
      routebox/.style={region, draw=teal!55!black, fill=teal!5},
      toolbox/.style={region, draw=orange!70!black, fill=orange!7},
      databox/.style={region, draw=cyan!55!black, fill=cyan!5},
      runtimebox/.style={region, draw=red!55!black, fill=red!4},
      qualitybox/.style={region, draw=yellow!55!black, fill=yellow!7},
      basebox/.style={region, draw=gray!70!black, fill=gray!7},
      demandcell/.style={cell, draw=blue!55!black, fill=blue!8, minimum width=2.82cm, minimum height=0.62cm},
      outcell/.style={cell, draw=magenta!55!black, fill=magenta!8, minimum height=0.42cm},
      llmcell/.style={cell, draw=green!50!black, fill=green!9, minimum height=0.42cm},
      routecell/.style={cell, draw=teal!55!black, fill=teal!9, minimum height=0.42cm},
      toolcell/.style={cell, draw=orange!70!black, fill=orange!11, minimum height=0.42cm},
      datacell/.style={cell, draw=cyan!55!black, fill=cyan!9, minimum height=0.42cm},
      govcell/.style={cell, draw=red!55!black, fill=red!7, minimum height=0.42cm},
      qualitycell/.style={cell, draw=yellow!55!black, fill=yellow!15, minimum width=2.82cm, minimum height=0.42cm},
      basecell/.style={cell, draw=gray!70!black, fill=gray!12, minimum width=1.34cm}
    ]
      \draw[roofbox, rounded corners=2pt]
        (-9.15, 6.16) -- (0, 7.66) -- (9.15, 6.16) -- cycle;
      \node[font=\bfseries\small] at (0, 6.98) {框架价值与意义};
      \node[note] at (0, 6.46) {从单图检测事实到项目级韧性评估闭环，形成可追溯、可复核、可扩展、可部署、可恢复的 Agent 工程方案};

      \node[demandbox, minimum width=3.20cm, minimum height=11.56cm] (demand) at (-7.55, -0.03) {};
      \node[sideTitle] at (-7.55, 5.42) {业务需求与问题输入};
      \node[demandcell] at (-7.55, 4.78) {项目级图像资产};
      \node[demandcell] at (-7.55, 3.68) {图像组 / 立面区域};
      \node[demandcell] at (-7.55, 2.59) {多材质幕墙图像};
      \node[demandcell] at (-7.55, 1.49) {批量上传与状态追踪};
      \node[demandcell] at (-7.55, 0.40) {单图单缺陷检测\\难以形成项目结论};
      \node[demandcell] at (-7.55, -0.70) {人工汇总成本高\\报告形成滞后};
      \node[demandcell] at (-7.55, -1.79) {用户真实目标\\建筑级韧性评估};
      \node[demandcell] at (-7.55, -2.89) {可复核证据\\维护优先级};
      \node[demandcell] at (-7.55, -3.98) {多区域风险\\分级处置};
      \node[demandcell] at (-7.55, -5.08) {报告成果\\归档交付};

      \node[outputbox, minimum width=10.05cm, minimum height=1.16cm] (output) at (0, 5.20) {};
      \node[layerTitle] at (-4.92, 5.50) {目标输出层};
      \node[outcell, minimum width=1.85cm] at (-3.425, 4.98) {单图检测总结};
      \node[outcell, minimum width=1.85cm] at (-1.142, 4.98) {项目评估报告};
      \node[outcell, minimum width=1.85cm] at (1.142, 4.98) {区域风险等级};
      \node[outcell, minimum width=1.85cm] at (3.425, 4.98) {维护改进建议};

      \node[llmbox, minimum width=10.05cm, minimum height=1.95cm] (llm) at (0, 3.33) {};
      \node[layerTitle] at (-4.92, 4.03) {LLM 分层信息整合层};
      \node[llmcell, minimum width=1.85cm] at (-3.425, 3.37) {原子事实层\\指标与产物};
      \node[llmcell, minimum width=1.85cm] at (-1.142, 3.37) {图像语义层\\图像级总结};
      \node[llmcell, minimum width=1.85cm] at (1.142, 3.37) {项目报告层\\多源融合};
      \node[llmcell, minimum width=1.85cm] at (3.425, 3.37) {报告结构层\\固定字段};
      \node[llmcell, minimum width=2.55cm] at (-3.075, 2.69) {结构化数据输入};
      \node[llmcell, minimum width=2.55cm] at (0.00, 2.69) {输出压缩 JSON 校验};
      \node[llmcell, minimum width=2.55cm] at (3.075, 2.69) {一致性约束};

      \node[routebox, minimum width=10.05cm, minimum height=1.72cm] (routing) at (0, 1.17) {};
      \node[layerTitle] at (-4.92, 1.75) {Agent 自主任务调度层};
      \node[routecell, minimum width=1.85cm] at (-3.425, 1.23) {幕墙场景识别};
      \node[routecell, minimum width=1.85cm] at (-1.142, 1.23) {材质语义识别};
      \node[routecell, minimum width=1.85cm] at (1.142, 1.23) {任务代码生成};
      \node[routecell, minimum width=1.85cm] at (3.425, 1.23) {非幕墙过滤};
      \node[routecell, minimum width=2.55cm] at (-3.075, 0.67) {枚举归一去重排序};
      \node[routecell, minimum width=2.55cm] at (0.00, 0.67) {子任务创建策略};
      \node[routecell, minimum width=2.55cm] at (3.075, 0.67) {失败回调重试入口};

      \node[toolbox, minimum width=10.05cm, minimum height=1.72cm] (tools) at (0, -0.87) {};
      \node[layerTitle] at (-4.92, -0.29) {原子检测工具层};
      \node[toolcell, minimum width=1.42cm] at (-3.64, -0.81) {金属锈蚀};
      \node[toolcell, minimum width=1.42cm] at (-1.82, -0.81) {石材裂缝};
      \node[toolcell, minimum width=1.42cm] at (0.00, -0.81) {石材污渍};
      \node[toolcell, minimum width=1.42cm] at (1.82, -0.81) {玻璃平整度};
      \node[toolcell, minimum width=1.42cm] at (3.64, -0.81) {玻璃爆裂};
      \node[toolcell, minimum width=1.85cm] at (-3.425, -1.37) {模型注册};
      \node[toolcell, minimum width=1.85cm] at (-1.142, -1.37) {运行目录隔离};
      \node[toolcell, minimum width=1.85cm] at (1.142, -1.37) {模型懒加载};
      \node[toolcell, minimum width=1.85cm] at (3.425, -1.37) {动态产物清单};

      \node[databox, minimum width=10.05cm, minimum height=1.72cm] (data) at (0, -2.91) {};
      \node[layerTitle] at (-4.92, -2.33) {数据证据与上下文层};
      \node[datacell, minimum width=1.85cm] at (-3.425, -2.85) {项目信息};
      \node[datacell, minimum width=1.85cm] at (-1.142, -2.85) {图像组};
      \node[datacell, minimum width=1.85cm] at (1.142, -2.85) {图像资产};
      \node[datacell, minimum width=1.85cm] at (3.425, -2.85) {检测产物};
      \node[datacell, minimum width=2.55cm] at (-3.075, -3.41) {区域指标};
      \node[datacell, minimum width=2.55cm] at (0.00, -3.41) {人工复核};
      \node[datacell, minimum width=2.55cm] at (3.075, -3.41) {报告源数据};

      \node[runtimebox, minimum width=10.05cm, minimum height=1.72cm] (runtime) at (0, -4.95) {};
      \node[layerTitle] at (-4.92, -4.37) {工程约束与运行治理层};
      \node[govcell, minimum width=1.42cm] at (-3.64, -4.89) {状态机};
      \node[govcell, minimum width=1.42cm] at (-1.82, -4.89) {协议回调};
      \node[govcell, minimum width=1.42cm] at (0.00, -4.89) {事务落库};
      \node[govcell, minimum width=1.42cm] at (1.82, -4.89) {对象授权};
      \node[govcell, minimum width=1.42cm] at (3.64, -4.89) {事件广播};
      \node[govcell, minimum width=1.85cm] at (-3.425, -5.45) {失败重试};
      \node[govcell, minimum width=1.85cm] at (-1.142, -5.45) {并发超时};
      \node[govcell, minimum width=1.85cm] at (1.142, -5.45) {日志追踪};
      \node[govcell, minimum width=1.85cm] at (3.425, -5.45) {CronJob 清理};

      \node[qualitybox, minimum width=3.20cm, minimum height=11.56cm] (quality) at (7.55, -0.03) {};
      \node[sideTitle] at (7.55, 5.42) {质量治理与反馈};
      \node[qualitycell] at (7.55, 4.78) {提示词规则\\字段白名单};
      \node[qualitycell] at (7.55, 3.68) {输出合法性校验\\JSON 对象};
      \node[qualitycell] at (7.55, 2.59) {任务代码校验\\非法值拦截};
      \node[qualitycell] at (7.55, 1.49) {产物完整性校验\\Uploaded=true};
      \node[qualitycell] at (7.55, 0.40) {报告一致性校验\\图像组与图像数};
      \node[qualitycell] at (7.55, -0.70) {人机协同复核\\准确 / 不准确 / 评论};
      \node[qualitycell] at (7.55, -1.79) {WebSocket 反馈\\进度与状态};
      \node[qualitycell] at (7.55, -2.89) {失败重试闭环\\检测与报告};
      \node[qualitycell] at (7.55, -3.98) {可扩展接入\\新增模型能力};
      \node[qualitycell] at (7.55, -5.08) {日志观测\\问题定位};

      \node[basebox, minimum width=18.30cm, minimum height=0.95cm] (base) at (0, -6.61) {};
      \node[layerTitle] at (-8.65, -6.61) {基础支撑};
      \node[basecell] at (-6.10, -6.61) {MySQL};
      \node[basecell] at (-4.08, -6.61) {Redis};
      \node[basecell] at (-2.07, -6.61) {MinIO};
      \node[basecell] at (-0.05, -6.61) {RocketMQ};
      \node[basecell] at (1.97, -6.61) {gRPC};
      \node[basecell] at (3.99, -6.61) {Docker};
      \node[basecell] at (6.00, -6.61) {Log};
      \node[basecell] at (8.02, -6.61) {Config};

      \draw[arrow] (-5.95, 1.17) -- node[note, above] {需求牵引} (-5.025, 1.17);
      \draw[arrow] (-5.95, -2.91) -- node[note, above] {项目资产} (-5.025, -2.91);
      \draw[arrow] (5.025, 3.33) -- node[note, above] {输出约束} (5.95, 3.33);
      \draw[arrow] (5.025, 1.17) -- node[note, above] {任务治理} (5.95, 1.17);
      \draw[arrow] (5.025, -0.87) -- node[note, above] {产物校验} (5.95, -0.87);
      \draw[arrow] (5.025, -2.91) -- node[note, above] {复核反馈} (5.95, -2.91);

      \draw[arrow] (0.00, -6.13) -- node[note, right] {工程底座} (0.00, -5.81);
      \draw[arrow] (0.00, -4.09) -- node[note, right] {状态约束} (0.00, -3.77);
      \draw[arrow] (0.00, -2.05) -- node[note, right] {工具事实} (0.00, -1.73);
      \draw[arrow] (0.00, -0.01) -- node[note, right] {工具候选} (0.00, 0.31);
      \draw[arrow] (0.00, 2.03) -- node[note, right] {结构化输入} (0.00, 2.36);
      \draw[arrow] (0.00, 4.31) -- node[note, right] {专业输出} (0.00, 4.62);

    \end{tikzpicture}
  }
  \caption{Agent 驱动的幕墙缺陷检测与韧性评估方法框架}\label{fig:design-agent-method-framework}
\end{figure}

\cref{fig:design-agent-method-framework} 是对本文方法框架的整体组织。可以先从左右两侧理解这张图：左侧对应业务需求与问题输入，包括项目级图像资产、图像组、立面区域、多材质幕墙图像、批量上传和状态追踪等对象，这些内容来自\cref{chap:introduction}中提到的单图检测割裂、人工汇总成本高和报告形成滞后等问题；右侧对应质量治理与反馈，用来约束 Agent 输出、任务代码、检测产物、报告一致性、人工复核、WebSocket 状态反馈、失败重试和日志观测等环节。中间部分则表示系统内部的处理主线，即从工程支撑开始，逐步形成检测事实、语义整合结果和用户最终看到的评估结论。

沿图中的箭头看，业务需求先进入 Agent 自主任务调度层，系统据此判断每张图像需要执行哪些检测；项目资产进入数据证据与上下文层后，被整理为后续 Agent 和报告生成可以使用的结构化上下文。底部的基础支撑和工程治理负责状态、协议、对象授权、事务落库和失败恢复。再往上，原子检测工具层生成可以追溯到图像和任务的检测事实，LLM 分层信息整合层在这些结构化输入基础上生成图像级总结和项目报告。目标输出层面向用户展示单图检测总结、项目评估报告、区域风险等级和维护改进建议。这样安排的目的，是把不确定的语义判断放在受控 Agent 节点中，把状态推进、数据管理和质量校验交给工程系统处理。

框架底部是基础支撑层，主要由 MySQL、Redis、MinIO、RocketMQ、gRPC、Docker、日志和配置等组成。这些设施本身并不产生智能结论，但它们决定了系统能否稳定保存状态、存储图像产物、传递事件、完成服务间通信，并支持容器化部署。其上一层是工程约束与运行治理，包括状态机、协议回调、事务落库、对象授权、事件广播、失败重试、并发超时、日志追踪和 CronJob 清理等机制。本文将这些机制单独放在一层，是因为 Agent 不能直接写库、直接推进项目状态，也不能绕过产物完整性校验；智能判断要进入业务流程，必须先经过这些工程边界。

数据证据与上下文层保存系统决策和报告生成所需的事实依据，例如项目信息、图像组、图像资产、检测产物、区域指标、人工复核和报告源数据。原本分散在单张图像、多个任务和多类缺陷中的信息，在这一层被组织为可以查询、复核和追溯的上下文。原子检测工具层则负责封装团队已有的检测模型和图像处理脚本，将金属锈蚀、石材裂缝、石材污渍、玻璃平整度和玻璃爆裂五类能力接入统一任务调度框架。为了降低不同模型之间的耦合，系统还需要配合模型注册、运行目录隔离、模型懒加载和动态产物清单等设计。

Agent 自主任务调度层处理的是“哪些图像需要执行哪些检测”。它不直接给出缺陷结论，而是根据幕墙场景识别和材质语义识别结果生成有限任务代码，再通过枚举归一、去重排序、子任务创建策略和失败回调重试入口，将语义判断转换为可执行的工程动作。LLM 分层信息整合层处理的是“检测结果如何进入项目级认知”。该层并不把所有底层指标一次性丢给模型，而是把原子事实层的指标与产物、图像语义层的图像级总结、项目报告层的多源融合和报告结构层的固定字段分开处理，并通过结构化数据输入、输出压缩 JSON 校验和一致性约束降低输出漂移与事实编造风险。

目标输出层最终面向用户提供单图检测总结、项目评估报告、区域风险等级和维护改进建议。右侧质量治理与反馈贯穿整个框架，主要用于检查提示词规则、字段白名单、JSON 对象合法性、任务代码合法性、产物完整性、报告一致性、人工复核、WebSocket 反馈、失败重试、扩展接入和日志观测等内容。由此可见，\cref{fig:design-agent-method-framework} 强调的是宏观方法组织，而不是某一次请求的技术细节。系统既不是让一个 Agent 从图像直接生成最终报告，也不是让所有模型无差别执行后再堆叠结果，而是把“证据生成”和“语义整合”分开，把“Agent 判断”和“业务状态推进”分开。分类 Agent 负责把图像语义映射为有限任务代码，Reasoning 服务负责生成可追溯的检测事实，图像总结 Agent 负责把底层指标转换为单图语义，项目报告 Agent 再结合项目背景、图像组和复核意见形成建筑级结论。

该方法与\cref{sec:foundation-agent-engineering-controllability-problem}讨论的 Agent 工程化可控性问题相对应。Agent 具有语义理解和自然语言组织优势，但直接赋予 Agent 自由行动权限会带来输出不稳定、工具选择错误和状态误推进等风险。因此，本文采用“动作空间有限、输出结构固定、业务状态外置”的设计。Agent 的输出必须经过协议解析和业务校验，所有数据库写入和项目状态推进均由 Core BIZ 执行。总体而言，该框架的核心在于：Agent 负责受控语义决策和分层信息整合，工程系统负责边界、状态、权限、数据、质量治理和异常恢复。

\section{Agent 在系统中的角色边界}
\label{sec:design-agent-role-definition}

在本系统中，Agent 不是一个开放式对话入口，也不是可以自主操作数据库的通用代理。本系统将 Agent 限定为业务流程中的三个受控节点：分类 Agent、图像总结 Agent 和项目报告 Agent。三者对应从“任务调度”到“局部解释”再到“整体评估”的递进关系。

\begin{table}[!htb]
  \caption{系统 Agent 节点设计}\label{tab:agent-node-design}
  \centering
  \begin{tabular}{@{}p{0.14\linewidth}p{0.26\linewidth}p{0.28\linewidth}p{0.22\linewidth}@{}} \toprule
    \textbf{Agent 节点} & \textbf{输入} & \textbf{输出} & \textbf{业务作用} \\ \midrule
    分类 Agent & 待检测幕墙图像 & 检测任务代码列表 & 判断该图像应执行哪些原子检测能力 \\
    图像总结 Agent & 单张图像的结构化检测结果 & 固定字段的单图检测总结 JSON & 将底层指标压缩为可读的图像级语义 \\
    项目报告 Agent & 项目基础信息、图像组、图像级总结、检测结果和复核意见 & 固定字段的项目评估报告 JSON & 形成建筑幕墙韧性评估结论和维护建议 \\ \bottomrule
  \end{tabular}
\end{table}

\cref{tab:agent-node-design} 表明，本文所说的 Agent 不是“大模型能力”的笼统称呼，而是具有固定业务入口、固定输入来源和固定输出契约的智能节点。分类 Agent 不输出缺陷结论，只输出任务代码；图像总结 Agent 不重新判断图像，只总结已完成检测结果；项目报告 Agent 不读取数据库，只依据系统组织好的项目上下文生成报告。这样的角色边界使 Agent 能够参与业务流程，但不会成为业务状态的所有者。

角色边界还体现在权限控制上。Agent 不能直接调用原子检测模型，不能直接访问 MySQL，不能直接写入 MinIO，也不能直接推进项目阶段。所有 Agent 输出都必须先返回 Activity 服务，再由 Activity 服务通过 gRPC 回调 Core BIZ。

从设计原则上看，本文采用的是“智能节点最小授权”思想。分类 Agent 只获得图像，不获得数据库写权限；总结 Agent 只获得结构化 JSON，不获得项目状态推进权限；项目报告 Agent 只获得系统整理后的多源检测结构化数据，不直接访问数据库和对象存储目录。这样可以降低 Agent 输出异常时的影响范围。即便某一次 Agent 输出不合规，最多导致当前节点失败，不会破坏项目主数据，也不会造成跨项目数据污染。

\section{技术路线取舍分析}
\label{sec:design-technical-route-tradeoff-analysis}

本小节解释为什么系统采用“受控 Agent + 原子检测能力 + 分层整合”的路线，而不是采用固定检测流程、端到端大模型或开放式 Agent 等更直观的方案。

\subsection{方案评价维度}
\label{subsec:design-scheme-evaluation-dimensions}

\begin{table}[!htb]
  \caption{技术方案评价维度}\label{tab:design-scheme-evaluation-dimensions}
  \centering
  \begin{tabular}{@{}p{0.16\linewidth}p{0.38\linewidth}p{0.38\linewidth}@{}} \toprule
    \textbf{评价维度} & \textbf{设计含义} & \textbf{本文关注点} \\ \midrule
    任务选择能力 & 系统能否根据图像材质和缺陷线索自动选择检测能力 & 避免人工逐图选择模型，也避免所有图像无差别执行全量检测 \\
    证据可追溯性 & 用户能否追溯结论来自哪张图像、哪个检测任务和哪些产物 & 原图、掩码、标注图、区域指标和人工复核都应被保留 \\
    状态可控性 & 长链路任务失败时能否定位失败节点并恢复 & 分类、子检测、产物上传、总结和报告生成必须有独立状态 \\
    报告可组织性 & 系统能否把大量单图结果转化为建筑级评估结论 & 需要支持图像组、区域差异、风险等级和维护建议 \\
    工程可扩展性 & 后续能否新增检测能力、展示字段或报告字段 & 任务代码、模型注册、结构化结果和前端展示应保持松耦合 \\
    用户可信性 & 输出是否符合工程用户对安全、责任和复核的要求 & Agent 结论不能覆盖事实，人工复核必须进入最终报告上下文 \\ \bottomrule
  \end{tabular}
\end{table}

建筑幕墙韧性评估不是单纯的图像分类任务。系统既要能处理多张、多组、多材质图像，又要让土木工程用户能够追溯缺陷证据、查看中间产物、补充人工意见并获得项目级报告。因此，技术方案不能只用模型精度或开发难度评价，还需要从业务角度建立评价维度。

如\cref{tab:design-scheme-evaluation-dimensions} 所示，本文关注的是一个工程系统的综合可用性，而不是单一模型调用能力。若方案只强调端到端生成文本，可能在演示时看起来简单，但无法满足证据追踪和人工复核要求；若方案只强调数据库和流程控制，则又难以解决材质识别、任务选择和报告语言组织问题。本文技术路线正是在这两类需求之间寻找平衡。

\subsection{技术路线比较}
\label{subsec:design-technical-route-comparison}

本文在方案设计阶段形成了六类可选技术路线。它们在特定场景下都具有一定合理性，但适用范围和风险不同。为突出选择过程，本文从“如何选择检测任务”和“如何生成项目报告”两个关键问题出发进行比较。

\begin{table}[!htb]
  \caption{不同技术路线适用性比较}\label{tab:design-technical-route-comparison}
  \centering
  \begin{tabular}{@{}p{0.16\linewidth}p{0.3\linewidth}p{0.3\linewidth}p{0.16\linewidth}@{}} \toprule
    \textbf{方案} & \textbf{优势} & \textbf{主要问题} & \textbf{本文取舍} \\ \midrule
    固定全量检测流程 & 实现简单，所有图像执行相同检测，状态容易控制 & 大量无关模型调用会浪费资源；非适用材质可能产生干扰性结果；用户难以理解为什么某类图像出现无关检测项 & 不采用 \\
    人工选择检测能力 & 人的判断准确性较高，便于结合现场经验 & 图像数量大时成本高、效率低；用户需要理解五类模型能力；难以支撑批量项目 & 作为复核补充 \\
    规则引擎路由 & 可解释性强，可根据文件夹、图像组或用户选项触发检测 & 规则依赖人工配置；难以处理混合材质图像；无法利用图像语义线索 & 不作为主方案 \\
    端到端大模型 & 链路短，理论上可以从图像或文本直接生成结论 & 缺少稳定区域指标和图像产物；上下文容易超限；失败节点不可定位；结论证据链薄弱 & 不作为主链路 \\
    开放式 Agent & 自主性强，理论上可自行规划和调用工具 & 工具边界不清晰，容易越权调用、状态误推进和输出不可控；不适合直接接管工程状态 & 不采用开放式权限 \\
    本系统方案 & 任务选择具备语义能力，像素级事实可追溯，报告生成分层可控，状态由工程系统统一推进 & 工程链路更复杂，需要协议、状态机、Worker、回调和产物管理协同实现 & 采用 \\ \bottomrule
  \end{tabular}
\end{table}

如\cref{tab:design-technical-route-comparison} 所示，本文没有选择端到端大模型，核心原因不是大模型能力不足，而是建筑幕墙评估需要“可追溯、可复核、可恢复”。土木工程用户不仅需要一段自然语言，还需要知道某张图像是否检测过、执行了哪些检测、检测区域在哪里、对应产物能否查看、人工复核如何影响最终报告。如果让大模型直接从图像生成报告，系统很难解释某一结论来自哪张图像、哪类缺陷和哪项指标，也难以在失败后定位是图像理解、缺陷识别还是报告生成环节出错。

固定全量检测流程也不适合本文场景。建筑幕墙图像可能是玻璃、石材、金属或多材质组合，若所有图像都执行五类检测，一方面增加推理成本，另一方面会产生大量对用户没有意义的“未检测到某类缺陷”的结论。对于工程用户而言，关键问题不是系统是否调用了足够多模型，而是系统是否在合适图像上调用了合适模型，并能解释对应检测结果。

规则引擎方案在传统业务系统中常见，例如根据图像组名称、用户选择或固定项目模板决定检测类型。但在无人机巡检或批量图像上传场景中，图像组名称并不能可靠表示图像材质，同一图像中也可能同时包含玻璃和金属构件。纯规则方案难以处理视觉语义不确定性。本文因此引入分类 Agent，而不是完全依赖人工配置。

\subsection{技术路线组合逻辑}
\label{subsec:design-technical-route-composition}

本系统最终采用组合式技术路线，其组合逻辑可以概括为三点。第一，视觉语义判断交给 Agent，但任务动作集合必须固定。Agent 可以判断图像中是否存在金属、石材或玻璃，但只能输出系统预定义的五类任务代码。第二，像素级事实交给原子检测能力。锈蚀区域、裂缝区域、污渍占比、平整度指标和爆裂置信度都来自专门模型或图像处理流程，而不是由 LLM 直接编写。第三，工程结论交给 LLM 分层整合。LLM 不负责发现像素级缺陷，而负责把结构化结果组织为用户可读、专业一致的评估语言。

这种组合使系统在智能性和可控性之间取得平衡。Agent 自主任务调度方法解决“做哪些检测”的语义选择问题；原子检测能力工具化解决“怎样获得像素级和区域级事实”的问题；LLM 分层信息整合策略解决“如何从局部事实形成工程评估结论”的问题；工程状态机解决“如何保证流程可追踪和可恢复”的问题。也就是说，Agent 并不是替代工程系统，而是补足传统工程系统在语义判断和信息整合上的不足。

该方案的代价是工程复杂度提高。系统需要维护任务代码、检测主任务、子任务、总结任务、报告任务、对象存储产物、人工复核和事件通知等多个对象。但这些复杂度是为了解决真实业务痛点而引入的必要复杂度，而不是形式上的服务拆分。通过将复杂度显式建模，系统才能支撑项目级韧性评估，而不是只完成单次图像推理。

\subsection{工业界实践启示}
\label{subsec:design-industry-practice-inspiration}

从工业界智能体系统实践看，可靠的 Agent 应用通常不会让大模型直接掌控全部业务状态，而是将大模型放在受限规划、信息抽取、内容生成或工具选择节点中。任务执行、权限控制、状态管理和数据持久化仍由传统工程系统完成。本文方案借鉴这一思路，将 Agent 定位为“受控决策节点”和“语义整合节点”，而不是完整业务流程的所有者。

这种设计尤其适合建筑幕墙韧性评估场景。幕墙检测结果具有安全属性，报告结论可能影响后续维护决策和工程责任，系统不能只追求自然语言输出的流畅性，还必须保留证据链、状态链和人工复核链。因此，本文不是简单地把大模型接入前端，而是围绕 Agent 输出契约、任务代码枚举、项目状态机、对象存储产物和人工复核机制构建完整方案。

\section{Agent 自主任务调度方法}
\label{sec:design-agent-autonomous-task-scheduling-method}

\begin{figure}[!htbp]
  \centering
  \resizebox{0.94\linewidth}{!}{%
    \begin{tikzpicture}[
      font=\scriptsize,
      >=stealth,
      process/.style={draw, rounded corners=4pt, line width=0.62pt, align=center, minimum width=10.8cm, fill=blue!4, draw=blue!55!black, inner sep=4pt},
      agent/.style={process, fill=green!6, draw=green!50!black},
      rule/.style={process, fill=orange!8, draw=orange!70!black},
      contract/.style={process, fill=purple!6, draw=purple!55!black},
      biz/.style={draw, rounded corners=4pt, line width=0.62pt, align=center, minimum width=4.95cm, minimum height=0.78cm, fill=cyan!6, draw=cyan!55!black},
      reasoning/.style={biz, fill=red!5, draw=red!55!black},
      arrow/.style={->, line width=0.62pt},
      note/.style={font=\tiny, align=center}
    ]
      \draw[draw=gray!70!black, fill=gray!7, rounded corners=4pt, line width=0.65pt]
        (-6.05, 0.82) rectangle (8.20, -9.52);

      \node[process, minimum height=0.76cm] (input) at (0, 0) {输入图像 $I$};

      \node[agent, minimum height=1.08cm] (semantic) at (0, -1.52) {
        \textbf{图像语义理解}\\[2pt]
        是否为建筑幕墙 \quad 幕墙材质识别 \quad 明显缺陷线索识别
      };

      \node[rule, minimum height=1.08cm] (candidate) at (0, -3.20) {
        \textbf{候选任务生成}\\[2pt]
        金属 $\rightarrow$ \texttt{corrosion}
        \quad 石材 / 混凝土 $\rightarrow$ \texttt{crack} + \texttt{stain}
        \quad 玻璃 $\rightarrow$ \texttt{flatness} + \texttt{spalling}
      };

      \node[contract, minimum height=1.08cm] (governance) at (0, -4.88) {
        \textbf{任务代码治理}\\[2pt]
        非幕墙图像过滤 \quad 非法任务代码拦截 \quad 去重 \quad 固定顺序排序 \quad 空集合允许
      };

      \node[process, minimum height=1.08cm, fill=teal!6, draw=teal!55!black] (set) at (0, -6.56) {
        \textbf{受控任务集合 $R(I)$}\\[2pt]
        $\{\texttt{corrosion},\texttt{crack},\texttt{stain},\texttt{flatness},\texttt{spalling}\}$ 的子集
      };

      \node[biz] (biz) at (-2.85, -8.69) {Core BIZ 创建检测子任务};
      \node[reasoning] (reasoning) at (2.85, -8.69) {Reasoning 执行原子检测};

      \draw[arrow] (input.south) -- (semantic.north);
      \draw[arrow] (semantic.south) -- (candidate.north);
      \draw[arrow] (candidate.south) -- (governance.north);
      \draw[arrow] (governance.south) -- (set.north);
      \draw[arrow] (set.south) -- (0, -7.70) -| (biz.north);
      \draw[arrow] (biz.east) -- (reasoning.west);
      \node[note, anchor=west] at (5.60, -1.52) {Agent 负责语义判断};
      \node[note, anchor=west] at (5.60, -4.88) {系统负责输出约束};
      \node[note, anchor=west] at (5.60, -8.69) {BIZ 负责状态推进};
    \end{tikzpicture}
  }
  \caption{Agent 自主任务调度方法流程}\label{fig:design-agent-autonomous-task-scheduling-flow}
\end{figure}

Agent 自主任务调度方法是本文核心方法之一。其目标是将“用户或工程师人工判断应该跑什么模型”的过程转化为受控的自动任务规划过程。具体而言，如\cref{fig:design-agent-autonomous-task-scheduling-flow} 所示，分类 Agent 接收单张图像，判断图像是否属于建筑幕墙场景，并识别其中可能存在的金属、石材/混凝土和玻璃等幕墙材质，再输出对应的原子检测任务代码。

\subsection{任务调度问题建模}
\label{subsec:design-task-scheduling-problem-modeling}

从形式上看，任务调度可以抽象为一个受限多标签选择问题。输入为项目中的单张图像，输出为任务代码集合：

\begin{equation}\label{eq:design-agent-task-routing-set}
  R(I) \subseteq \{\texttt{corrosion}, \texttt{crack}, \texttt{stain}, \texttt{flatness}, \texttt{spalling}\}
\end{equation}

其中，$I$ 表示图像，$R(I)$ 表示该图像需要执行的原子检测能力集合。与普通多标签分类不同，本文任务调度结果并不是最终用户结论，而是后续工程状态机的输入。因此，输出集合必须满足三个条件：第一，元素必须属于固定任务代码枚举；第二，集合允许为空，表示图像不适合或不需要检测；第三，集合中的每个任务代码最多出现一次，并应按统一顺序归一化，以便数据库保存、前端展示和日志排查。

这种问题建模的意义在于把 Agent 的开放语义理解能力压缩为工程系统可以消费的离散动作。对于 Agent 而言，判断图像中是否存在玻璃、石材、金属或明显缺陷线索属于语义判断；对于工程系统而言，真正可执行的动作只有五类原子检测能力。Agent 自主任务调度方法正是连接这两者的桥梁。

\subsection{材质语义与任务代码映射}
\label{subsec:design-material-to-task-code-mapping}

幕墙常见材质与检测任务之间存在较明确的对应关系。金属构件主要关注锈蚀问题；石材或混凝土面板主要关注裂缝和污渍问题；玻璃幕墙主要关注平整度异常和爆裂问题。本文将这一工程经验固化为任务映射规则，并通过提示词约束 Agent 的输出。

\begin{table}[!htb]
  \caption{幕墙材质语义与检测任务代码映射关系}\label{tab:design-material-task-code-mapping}
  \centering
  \begin{tabular}{@{}p{0.2\linewidth}p{0.2\linewidth}p{0.54\linewidth}@{}} \toprule
    \textbf{图像语义线索} & \textbf{对应任务代码} & \textbf{设计说明} \\ \midrule
    金属幕墙、金属连接件、明显锈蚀痕迹 & \texttt{corrosion} & 金属锈蚀会影响构件耐久性和连接可靠性，是金属类图像的重点检测能力 \\
    石材、混凝土、类似石材面板 & \texttt{crack}、\texttt{stain} & 石材类图像同时需要关注裂缝和污渍，前者偏安全与耐久性，后者偏外观和渗水线索 \\
    玻璃幕墙、玻璃反射面、明显玻璃破损 & \texttt{flatness}、\texttt{spalling} & 玻璃类图像既需要判断表面或安装状态异常，也需要判断是否存在爆裂风险 \\
    多材质组合幕墙 & 多类任务代码合并去重 & 对包含玻璃、金属和石材的图像，允许输出多个任务代码，后续并行执行 \\
    非建筑幕墙、无法判断、无检测意义图像 & 空数组 & 系统不强行执行检测，避免产生无意义结果和错误报告依据 \\ \bottomrule
  \end{tabular}
\end{table}

\cref{tab:design-material-task-code-mapping} 中的映射并不是简单关键词规则，而是提示 Agent 结合图像内容进行判断。比如，图像中只要出现大面积玻璃面板，即使没有明显爆裂，也应输出 \texttt{flatness} 和 \texttt{spalling}，因为检测任务的目的不仅是确认已知缺陷，还包括排查潜在异常；图像中若出现石材或混凝土表面，即使裂缝不明显，也需要执行 \texttt{crack} 和 \texttt{stain}，以获得可追溯检测结果。

\subsection{输出契约与服务端归一化}
\label{subsec:design-classification-output-contract}

在分类路由环节，系统只接收一种输出形式：单行压缩 JSON 字符串，并且对象中只保留 \texttt{task\_codes} 字段。该字段的值为任务代码字符串数组。例如，玻璃幕墙图像可以输出 \texttt{\{"task\_codes":["flatness","spalling"]\}}，非幕墙图像则输出 \texttt{\{"task\_codes":[]\}}。这样设计不是为了追求格式上的简洁，而是为了让后续 Worker 能拿到确定的任务输入。如果 Agent 额外输出自然语言解释、Markdown 代码块、中文任务名或其他字段，服务端就需要再猜测模型意图，子任务创建也会变得不稳定。

因此，服务端不会默认提示词约束一定生效，而是会对 Agent 返回内容再做一次归一化处理。处理过程包括去除多余空白、解析 JSON 对象、检查字段是否存在、检查 \texttt{task\_codes} 是否为字符串数组、过滤或拒绝非法任务代码，并对合法任务代码去重后按统一顺序排序。归一化后的结果再转换为系统枚举，进入 Core BIZ 的任务创建逻辑。这个过程相当于把模型的开放输出收回到工程系统可识别的有限集合中：Agent 负责给出判断，服务端负责校验、兜底和保护业务流程。

\subsection{任务调度闭环设计}
\label{subsec:design-task-scheduling-closed-loop}

Agent 自主任务调度并不是调用一次 Classification 服务就结束，而是检测主任务状态机中的前置环节。用户在前端启动 Agent 智能检测后，Core BIZ 会先为每张图像创建检测主任务；后台 Worker 再把主任务推进到分类中状态，获取图像访问地址，并调用 Classification 服务。分类 Agent 返回任务代码后，Classification 服务通过回调把结果交回 Core BIZ。随后 Core BIZ 在事务中写入任务代码、创建对应检测子任务，并把主任务推进到子任务检测阶段。

这里需要特别区分两个边界。一方面，分类 Agent 只回答“应该做什么检测”，它不负责判断主任务是否完成，也不直接创建数据库子任务。是否进入检测阶段、是否发布状态事件，都由 Core BIZ 按照事务结果统一处理。另一方面，调度结果需要和后续检测结果绑定。主任务记录会保存本次生成的任务代码，用户以后查看结果时，可以知道系统为什么执行了这些子任务；项目报告在整合结论时，也可以据此判断某类检测是否真正执行过。

\subsection{异常场景与边界处理}
\label{subsec:design-routing-exception-handling}

任务调度中存在多种异常情况。第一，Agent 输出空任务列表。该情况并不一定是失败，可能表示图像不是建筑幕墙、图像质量不足、或无法可靠判断材质。系统应将其视为一种有效结果，并根据业务规则结束当前图像检测链路。第二，Agent 输出格式错误。该情况属于分类节点失败，应回调失败状态，允许用户后续重试。第三，Agent 输出合法但不符合实际材质。该情况需要通过人工复核和项目报告上下文进行纠偏，而不是在系统层面强行修改历史检测事实。

本文没有采用“输出不确定时执行全部检测”的策略。原因在于，全部检测虽然看似保险，但会产生大量非适用检测结果，反而降低报告可信度。对于无法可靠判断的图像，系统更适合输出空任务，并由用户根据图像质量或项目需要重新上传、重新分组或人工补充说明。这样的策略符合土木工程场景对结论依据的审慎要求。

\section{原子检测能力工具化}
\label{sec:design-atomic-detection-capability-tooling}

原子检测能力工具化是指将团队前期积累的检测模型和图像处理脚本封装为统一可调度能力。\cref{sec:foundation-curtain-wall-defect-detection-technology}已经介绍五类原子检测能力的技术基础，本节关注它们如何在本系统中成为 Agent 可间接调度的工具。如\cref{fig:design-atomic-detection-tooling-flow} 所示，工具化设计的关键不是重新提出某个检测模型，而是让不同来源、不同输入输出形式、不同产物数量的检测能力在同一工程协议下运行。

\begin{figure}[!htb]
  \centering
  \resizebox{0.94\linewidth}{!}{%
    \begin{tikzpicture}[
      font=\scriptsize,
      >=stealth,
      process/.style={draw, rounded corners=4pt, line width=0.62pt, align=center, minimum width=10.8cm, minimum height=1.08cm, fill=blue!4, draw=blue!55!black, inner sep=4pt},
      registry/.style={process, fill=green!6, draw=green!50!black},
      runtime/.style={process, fill=orange!8, draw=orange!70!black},
      model/.style={process, fill=purple!6, draw=purple!55!black},
      artifact/.style={process, fill=teal!6, draw=teal!55!black},
      result/.style={process, fill=cyan!6, draw=cyan!55!black},
      callback/.style={process, fill=red!5, draw=red!55!black},
      arrow/.style={->, line width=0.62pt},
      note/.style={font=\tiny, align=center}
    ]
      \draw[draw=gray!70!black, fill=gray!7, rounded corners=4pt, line width=0.65pt]
        (-6.05, 0.82) rectangle (8.20, -10.90);

      \node[process, minimum height=0.76cm] (task) at (0, 0) {检测子任务};

      \node[registry] (registry) at (0, -1.52) {
        \textbf{工具选择与模型注册}\\[2pt]
        \texttt{task\_code} 匹配检测能力 \quad 模型路径与运行参数加载 \quad 动态产物清单生成
      };

      \node[runtime] (workdir) at (0, -3.20) {
        \textbf{运行目录隔离}\\[2pt]
        创建子任务工作目录 \quad 下载原图 / 区域图 \quad 初始化产物目录
      };

      \node[model] (execute) at (0, -4.88) {
        \textbf{原子检测执行}\\[2pt]
        金属锈蚀 \quad 石材裂缝 \quad 石材污渍 \quad 玻璃平整度 \quad 玻璃爆裂
      };

      \node[artifact] (upload) at (0, -6.56) {
        \textbf{产物上传与完整性校验}\\[2pt]
        产物上传 MinIO \quad \texttt{Uploaded=true} 校验 \quad 产物数量与动态清单一致
      };

      \node[result] (result) at (0, -8.24) {
        \textbf{结构化结果生成}\\[2pt]
        缺陷状态 \quad 区域指标 \quad 图像产物 Key \quad 检测 Summary 原始数据
      };

      \node[callback] (callback) at (0, -9.92) {
        \textbf{Reasoning 回调 Core BIZ}\\[2pt]
        成功：写入子任务结果 \quad 失败：记录失败原因并等待重试
      };

      \draw[arrow] (task.south) -- (registry.north);
      \draw[arrow] (registry.south) -- (workdir.north);
      \draw[arrow] (workdir.south) -- (execute.north);
      \draw[arrow] (execute.south) -- (upload.north);
      \draw[arrow] (upload.south) -- (result.north);
      \draw[arrow] (result.south) -- (callback.north);

      \node[note, anchor=west] at (5.60, -1.52) {注册负责能力发现};
      \node[note, anchor=west] at (5.60, -3.20) {目录负责并发隔离};
      \node[note, anchor=west] at (5.60, -6.56) {校验负责结果可信};
      \node[note, anchor=west] at (5.60, -9.92) {回调负责状态闭环};
    \end{tikzpicture}
  }
  \caption{原子检测能力工具化执行流程}\label{fig:design-atomic-detection-tooling-flow}
\end{figure}

\subsection{统一注册机制}
\label{subsec:design-atomic-capability-registration}

工具化的第一层是任务代码统一。五类能力通过统一任务代码注册，Reasoning 服务启动时读取模型目录，并检查每个任务代码下的模型文件是否存在、是否为目录、是否为 Git LFS 指针文件。这样可以在服务启动阶段发现模型缺失或部署错误，而不是等到用户启动检测后才失败。

注册信息不仅包括任务代码，还包括模型名称、Python 执行入口和结果解析方式。通过注册机制，工程系统可以把不同模型统一看作“输入原图、输出报告和产物”的可调用工具。新增检测能力时，不需要改变 Agent 调用方式和 Worker 编排框架，只需新增任务代码、注册元数据、模型目录和 Python 检测入口。

\subsection{统一执行协议}
\label{subsec:design-atomic-execution-protocol}

工具化的第二层是执行协议统一。无论内部模型使用 YOLO、ResNet、OpenCV 还是其他图像处理流程，对外都遵守相同执行约定：输入为项目图像原图，运行目录按任务代码和图像 UUID 隔离，输出为 \texttt{report.json} 和若干图像产物。Golang 层读取并压缩报告 JSON，上传图像产物，再将结构化结果和产物清单回调给 Core BIZ。Python 模型不直接访问数据库，也不直接参与业务状态推进。

运行目录隔离是工具化中的重要设计。每个子任务拥有独立目录，目录中包含输入图像、模型生成的中间产物和报告文件。这样可以避免并发任务之间相互覆盖，也便于失败后定位具体任务。对于需要重试的任务，系统可以重新创建运行目录并重新执行，而不依赖上一次执行遗留文件。

通过统一执行协议，原本异构的检测模型和图像处理脚本被封装为可调度、可复核、可重试和可恢复的工程工具，为上层 Agent 调度和后续信息整合提供稳定的事实生成接口。

\subsection{统一结果契约}
\label{subsec:design-atomic-result-contract}

工具化的第三层是结果契约统一。不同任务的指标不同，例如锈蚀检测关注锈蚀区域数量、置信度、像素数和占比，裂缝检测关注裂缝像素与占比，污渍检测关注污渍区域和污渍占比，平整度检测关注边缘、梯度、直线和频谱等指标，爆裂检测关注是否爆裂与置信度。但这些结果都必须被整理为可落库、可展示、可总结的结构化字段。只有具备这种契约，原子检测能力才能从实验脚本变成工程系统中的工具。

\begin{table}[!htb]
  \caption{五类原子检测能力的主要结构化指标与图像产物}\label{tab:design-atomic-capability-output}
  \centering
  \begin{tabular}{@{}p{0.1\linewidth}p{0.4\linewidth}p{0.4\linewidth}@{}} \toprule
    \textbf{检测能力} & \textbf{主要结构化指标} & \textbf{主要图像产物} \\ \midrule
    金属锈蚀 & 是否锈蚀、区域数量、平均置信度、最高置信度、锈蚀像素数、锈蚀占比、区域列表 & 标注图、区域框、必要的检测可视化结果 \\
    石材裂缝 & 是否存在裂缝、裂缝区域数量、裂缝像素数、裂缝占比、区域列表 & 裂缝浮层图、裂缝图、区域可视化结果 \\
    石材污渍 & 是否存在污渍、污渍区域数量、平均污渍占比、最高污渍占比、区域列表 & 总体标注图、区域裁剪图、区域污渍结果图 \\
    玻璃平整度 & 结果类型、异常区域数量、边缘异常、梯度异常、直线异常、频谱异常和相关数值指标 & 区域原图、梯度图、直线图、频谱图 \\
    玻璃爆裂 & 是否爆裂、置信度、分类结果 & 原图或分类结果相关报告产物 \\ \bottomrule
  \end{tabular}
\end{table}

\cref{tab:design-atomic-capability-output} 展示了工具化输出的差异。系统不要求所有检测能力输出完全相同字段，而是要求它们遵守共同的外层契约：是否成功、报告 JSON、产物清单、区域指标和任务状态。这样既保留每类检测模型的专业性，又保证上层 Agent 总结和前端展示能够使用统一入口读取结果。

\subsection{工具化与 Agent 调度的关系}
\label{subsec:design-tooling-and-agent-routing}

原子检测能力并不直接暴露给用户，也不直接由 Agent 调用。Agent 输出的是任务代码，Core BIZ 根据任务代码创建子任务，Reasoning 服务再根据子任务代码查找已注册检测能力。这样的间接调用方式使系统能够控制执行顺序、并发、权限、失败重试和结果落库。若让 Agent 直接调用 Python 工具，则难以保证数据库状态一致，也难以在产物上传失败时阻止任务被误判为成功。

因此，本文中的“工具”不是大模型工具调用框架中的自由工具，而是工程系统管理下的受控工具。Agent 决定选择哪些工具，工程系统决定何时执行、如何执行、结果是否有效以及是否推进状态。这一点是本系统方案区别于开放式 Agent 的核心。

\section{LLM 分层信息整合策略}
\label{sec:design-large-language-model-hierarchical-information-integration-strategy}

LLM 分层信息整合策略是本文从单图检测走向项目级韧性评估的关键设计。旧式做法通常把每个模型的检测结果直接展示给用户，用户再人工汇总；另一种极端做法是把所有检测结果一次性输入大模型并要求其生成报告。前者无法解决报告滞后问题，后者容易受到上下文长度、低层指标噪声和输出不稳定影响。如\cref{fig:design-llm-hierarchical-integration-flow} 所示，本文采用分层整合方案，将信息处理划分为原子事实层、图像语义层和项目报告层。

\begin{figure}[!htbp]
  \centering
  \resizebox{0.94\linewidth}{!}{%
    \begin{tikzpicture}[
      font=\scriptsize,
      >=stealth,
      process/.style={draw, rounded corners=4pt, line width=0.62pt, align=center, minimum width=10.8cm, minimum height=1.08cm, fill=blue!4, draw=blue!55!black, inner sep=4pt},
      fact/.style={process, fill=green!6, draw=green!50!black},
      imageSummary/.style={process, fill=orange!8, draw=orange!70!black},
      source/.style={process, fill=purple!6, draw=purple!55!black},
      projectAgent/.style={process, fill=teal!6, draw=teal!55!black},
      report/.style={process, fill=cyan!6, draw=cyan!55!black},
      validation/.style={process, fill=red!5, draw=red!55!black},
      arrow/.style={->, line width=0.62pt},
      note/.style={font=\tiny, align=center}
    ]
      \draw[draw=gray!70!black, fill=gray!7, rounded corners=4pt, line width=0.65pt]
        (-6.05, 0.82) rectangle (8.20, -10.90);

      \node[process, minimum height=0.76cm] (input) at (0, 0) {检测事实与项目上下文};

      \node[fact] (fact) at (0, -1.52) {
        \textbf{原子事实层}\\[2pt]
        缺陷状态 \quad 区域指标 \quad 置信度 / 占比 \quad 图像产物 Key
      };

      \node[imageSummary] (image) at (0, -3.20) {
        \textbf{图像语义层}\\[2pt]
        图像级总结 Agent 生成总体概述、五类缺陷总结和后续建议
      };

      \node[source] (source) at (0, -4.88) {
        \textbf{项目报告源数据层}\\[2pt]
        项目信息 \quad 图像组 \quad 检测结果 \quad 图像总结 \quad 人工复核 \quad 统计字段
      };

      \node[projectAgent] (agent) at (0, -6.56) {
        \textbf{附件输入与项目报告 Agent}\\[2pt]
        读取多源检测结构化数据，按图像组、缺陷类型、复核意见进行多源融合
      };

      \node[report] (report) at (0, -8.24) {
        \textbf{结构化项目评估报告}\\[2pt]
        项目背景 \quad 总体结论 \quad 区域评估 \quad 风险等级 \quad 韧性评估 \quad 维护建议
      };

      \node[validation] (validation) at (0, -9.92) {
        \textbf{输出压缩与一致性校验}\\[2pt]
        单行 JSON 对象 \quad 图像组与图像数校验 \quad 防止编造 \quad 前端字段化展示
      };

      \draw[arrow] (input.south) -- (fact.north);
      \draw[arrow] (fact.south) -- (image.north);
      \draw[arrow] (image.south) -- (source.north);
      \draw[arrow] (source.south) -- (agent.north);
      \draw[arrow] (agent.south) -- (report.north);
      \draw[arrow] (report.south) -- (validation.north);

      \node[note, anchor=west] at (5.60, -1.52) {事实负责可追溯};
      \node[note, anchor=west] at (5.60, -3.20) {单图负责语义压缩};
      \node[note, anchor=west] at (5.60, -6.56) {项目负责多源融合};
      \node[note, anchor=west] at (5.60, -9.92) {校验负责输出稳定};
    \end{tikzpicture}
  }
  \caption{LLM 分层信息整合策略流程}\label{fig:design-llm-hierarchical-integration-flow}
\end{figure}

\subsection{信息分层整合的必要性}
\label{subsec:design-llm-hierarchy-necessity}

项目级幕墙评估与普通文本摘要不同。输入数据既包括项目基础信息，又包括图像组、图像元数据、五类检测指标、区域列表、图像产物、图像级总结和人工复核意见。如果将所有底层 JSON 直接输入项目报告 Agent，一方面会造成上下文过长，另一方面低层细节会干扰模型对项目级风险的组织。相反，如果只把图像级总结输入项目报告 Agent，又会损失必要的结构化指标和复核依据。

本文因此采用信息分层整合策略。原子检测能力首先产生可追溯事实，图像总结 Agent 将每张图像的事实压缩为语义概要，项目报告 Agent 再同时读取项目结构、检测事实、图像总结和人工复核意见等。这样既避免项目报告直接处理过多像素级或区域级细节，又保留必要的证据字段。

\begin{table}[!htb]
  \caption{LLM 信息整合分层输入内容和输出作用}\label{tab:design-llm-hierarchical-integration}
  \centering
  \begin{tabular}{@{}p{0.12\linewidth}p{0.42\linewidth}p{0.4\linewidth}@{}} \toprule
    \textbf{层次} & \textbf{输入内容} & \textbf{输出作用} \\ \midrule
    原子事实层 & 五类原子检测能力输出的结构化结果，包括缺陷存在性、区域指标、置信度、占比和产物 & 保留可追溯、可展示、可复核的底层证据 \\
    图像语义层 & 单张图像实际执行过的检测结果 JSON & 生成固定字段的图像总结，说明该图像是否存在缺陷、主要问题和维护建议 \\
    项目报告层 & 项目基础信息、图像组、每张图像检测结果、图像总结和人工复核意见 & 生成项目级韧性评估报告，给出整体状态、区域差异、风险等级和维护建议 \\ \bottomrule
  \end{tabular}
\end{table}

\cref{tab:design-llm-hierarchical-integration} 表明，图像级总结和项目级报告不是两个重复的摘要任务，而是两个不同抽象层次。图像级总结只回答“这一张图像说明什么”，输出字段限定为总体概述、五类缺陷总结和后续建议。项目级报告回答“整个项目说明什么”，需要综合图像组、缺陷分布、人工复核和项目背景，形成项目背景、数据集概述、总体结论、区域评估、风险等级、韧性评估、维护建议和限制说明等内容。

\subsection{图像级总结设计}
\label{subsec:design-image-level-summary}

图像级总结用于把单张图像已经产生的多类检测结果整理成前端可读的结构化文本。它的输入来自已落库的检测结果，不重新执行模型推理，也不重新判断图像是否存在某类缺陷。输出字段包括总体概述、金属锈蚀、石材裂缝、石材污渍、玻璃平整度、玻璃爆裂和后续建议。前端展示时可以根据字段是否为空决定显示哪些标签，避免把没有证据支撑的内容展示给用户。

这一设计中特别需要区分“未执行检测”和“未发现缺陷”。例如，分类 Agent 没有选择金属锈蚀检测时，图像级总结不应写“未发现锈蚀”，因为系统实际并没有对该图像执行锈蚀检测。对于已经执行但未提示明显异常的任务，可以在总体概述中说明已执行任务未发现明显异常，但对应缺陷字段仍保持为空。这样处理后，图像级总结不会把检测范围之外的内容写成结论。

从生成控制角度看，图像级总结需要有明确的证据边界。自然语言生成模型容易为了让文字完整而补充未执行任务的描述，这在工程报告中会造成误导。本文通过提示词规则和字段设计限制这种情况：只有存在对应材质或缺陷，并且检测结果能够支撑时，才填写对应字段；如果没有发现缺陷，后续建议可以为空；如果某类任务未执行，则不能把未执行任务等同于无缺陷结论。

\subsection{项目级报告源数据设计}
\label{subsec:design-project-report-source-data}

项目报告生成设计的关键不在于让 Agent 写长文本，而在于如何把多源数据组织成可供 Agent 稳定阅读的上下文。Core BIZ 会先将项目基础信息、图像组、图像检测结果、图像级总结和复核意见组装为结构化 JSON，并以 \texttt{source.txt} 文件形式存入对象存储。项目报告 Agent 通过预签名 URL 下载该文件并作为附件读取。

采用 \texttt{source.txt} 附件有两个原因。第一，项目级数据可能较长，直接把压缩 JSON 放在消息正文中容易触发上下文限制，也不利于日志阅读。附件方式可以让 Agent 明确读取完整源数据。第二，报告源数据本身是一个重要中间产物，保存在对象存储中可以追溯报告依据。若项目报告存在争议，可以查看当时传给 Agent 的源数据，而不是只查看最终报告文本。

\begin{table}[!htb]
  \caption{项目级报告源数据设计}\label{tab:design-project-report-source-data}
  \centering
  \begin{tabular}{@{}p{0.2\linewidth}p{0.7\linewidth}@{}} \toprule
    \textbf{数据块} & \textbf{设计作用} \\ \midrule
    项目基础信息 & 提供项目名称、建筑名称、地址、建成年份、描述等报告背景信息 \\
    图像组与图像列表 & 保留项目图像组织结构，使报告能够按立面、区域或图像组描述缺陷分布 \\
    原子检测结果 & 提供五类检测能力的结构化事实，包括缺陷存在性、数值指标和区域列表 \\
    图像级总结 & 将单图多任务结果压缩为可读语义，帮助项目报告 Agent 更好理解单图重点 \\
    人工复核意见 & 将用户对 Agent 结果的准确性判断和补充评论纳入最终报告依据 \\
    统计字段 & 包括图像组数量和图像数量，用于防止报告 Agent 编造数据规模 \\ \bottomrule
  \end{tabular}
\end{table}

\cref{tab:design-project-report-source-data} 主要说明报告源数据在进入 Agent 之前需要做一次筛选。系统不会把数据库中的所有字段原样传入模型，例如内部 ID、任务主键、产物 SHA256、运行耗时以及部分图像坐标，这些字段更多用于系统内部追踪或产物校验，对项目报告本身帮助不大，因而会在组装上下文时被排除。相反，缺陷指标、区域数量、图像组结构、人工复核意见和图像级总结等信息会被保留下来，因为这些内容直接关系到报告中的风险判断和维护建议。这样处理后，报告源数据不会被底层实现字段占满，同时仍能保留形成建筑级韧性评估结论所需要的依据。

\subsection{项目级报告输出契约}
\label{subsec:design-project-report-output-contract}

项目级报告输出必须是结构化 JSON，而不是普通 Markdown 文本。原因在于前端需要按照字段展示项目背景、韧性评估目标、数据集概述、总体评估结论、总体风险等级、局部评估结论、局部风险等级、总体韧性评估、维护和改进建议以及限制说明等信息。结构化输出可以避免前端解析自然语言段落，也便于后续扩展报告打印、导出和字段化展示。

项目报告 Agent 的输出还需要满足一致性约束。报告中的图像组数量和图像数量必须来自 \texttt{source.txt}，不得自行估算；区域评估应基于源数据中的图像组组织，而不是编造不存在的立面；风险等级需要与检测事实和人工复核保持一致；限制说明需要指出图像检测无法完全替代现场复核。系统在接收回调时会将 Agent 输出压缩为单行 JSON，并校验其为合法 JSON 对象。

\subsection{幻觉与上下文漂移防止机制}
\label{subsec:design-llm-hallucination-control}

LLM 分层信息整合策略除了生成报告文本之外，还需要尽量控制幻觉和上下文漂移。项目级报告和普通摘要不同，报告中会出现项目名称、图像数量、立面区域、缺陷分布和风险判断等内容，这些信息一旦被模型自行补充，就会直接影响用户对建筑状态的理解。实际设计时，最需要避免的情况不是语言表述不够流畅，而是报告中出现源数据里没有的项目背景、图像组数量或风险结论。

针对这一问题，本文从输入、提示词和回调校验三个位置进行约束。输入层面，系统在报告源数据中显式写入 \texttt{group\_count} 和 \texttt{image\_count} 等统计字段，避免模型自行根据文本描述推测数量。提示词层面，要求报告内容必须基于附件中的 JSON，不允许补充不存在的项目背景或检测事实。回调层面，系统会对报告 JSON 进行压缩和合法性校验；如果必要字段缺失、字段类型不符合要求，或者整体结构无法解析，则将本次报告任务视为失败，而不是把异常文本继续写入正式报告结果。

这些约束不能从理论上完全消除大模型幻觉，但可以降低幻觉进入系统结果的概率。更重要的是，系统会将报告源数据保存为 \texttt{source.txt}，使报告依据能够被回看。用户如果认为报告内容不准确，可以把报告结论与源数据逐项对照，进一步判断问题来自检测结果、人工复核、源数据组织，还是 Agent 生成过程本身。对于工程评估场景而言，可追溯性比单纯追求语言完整度更重要，因为最终报告需要能够说明“结论依据来自哪里”。

\subsection{信息分层整合与工程认知路径}
\label{subsec:design-engineering-cognition-path}

LLM 分层信息整合策略体现了工程报告生成的认知路径。实际工程师撰写项目报告时，也不会直接从每个像素坐标开始写，而是先理解每张图像的主要现象，再比较不同区域之间的缺陷分布，最后给出整体风险和维护优先级。LLM 分层信息整合策略把这一过程固化为软件流程，使 Agent 的语言能力服务于工程语义组织，而不是替代检测事实。

从用户体验角度看，图像级总结和项目级报告也解决了不同层次用户需求。项目操作人员在检测阶段需要快速知道某张图是否存在异常，因此需要图像级总结；项目负责人或土木工程团队更关注建筑整体状态、区域风险和维护优先级，因此需要项目级报告。分层输出使同一套检测事实能够服务不同层次的业务阅读场景。

\section{任务状态流转与人机协同机制}
\label{sec:design-detection-task-state-transition}

\begin{table}[!htb]
  \caption{检测主任务状态枚举含义}\label{tab:detection-state-meaning}
  \centering
  \begin{tabular}{@{}p{0.18\linewidth}p{0.64\linewidth}@{}} \toprule
    \textbf{状态枚举} & \textbf{业务含义} \\ \midrule
    \texttt{pending} & 主任务已创建，等待后台 Worker 消费 \\
    \texttt{classifying} & 分类 Agent 正在判断该图像需要执行哪些原子检测能力 \\
    \texttt{detecting} & 原子检测子任务已创建，系统正在等待检测结果和产物上传 \\
    \texttt{summarizing} & 子任务结果已满足总结条件，图像总结 Agent 正在生成图像级总结 \\
    \texttt{succeeded} & 该图像检测链路已成功结束 \\
    \texttt{failed} & 必要节点失败，当前图像检测链路失败，可通过重试重新创建任务 \\ \bottomrule
  \end{tabular}
\end{table}

如\cref{fig:design-task-state-human-review-flow} 和\cref{tab:detection-state-meaning} 所示，任务状态流转是本文方法落地的骨架。每张图像对应一个检测主任务，主任务从 \texttt{pending} 开始，后台 Worker 消费后进入 \texttt{classifying}。分类成功后，如果 Agent 输出了需要执行的任务代码，系统创建对应子任务并进入 \texttt{detecting}；如果 Agent 输出空任务列表，说明当前图像不需要进入原子检测链路，主任务可以直接成功结束。子任务全部成功后，系统创建图像总结任务并进入 \texttt{summarizing}；图像总结成功后，主任务进入 \texttt{succeeded}；任一必要节点失败，主任务进入 \texttt{failed}。

\begin{figure}[!htbp]
  \centering
  \resizebox{0.94\linewidth}{!}{%
    \begin{tikzpicture}[
      font=\scriptsize,
      >=stealth,
      process/.style={draw, rounded corners=4pt, line width=0.62pt, align=center, minimum width=10.8cm, minimum height=1.08cm, fill=blue!4, draw=blue!55!black, inner sep=4pt},
      classify/.style={process, fill=green!6, draw=green!50!black},
      detect/.style={process, fill=orange!8, draw=orange!70!black},
      summary/.style={process, fill=purple!6, draw=purple!55!black},
      state/.style={process, fill=teal!6, draw=teal!55!black},
      review/.style={process, fill=cyan!6, draw=cyan!55!black},
      report/.style={process, fill=red!5, draw=red!55!black},
      arrow/.style={->, line width=0.62pt},
      note/.style={font=\tiny, align=center}
    ]
      \draw[draw=gray!70!black, fill=gray!7, rounded corners=4pt, line width=0.65pt]
        (-6.05, 0.82) rectangle (8.20, -10.90);

      \node[process, minimum height=0.76cm] (start) at (0, 0) {检测主任务创建};

      \node[classify] (classifying) at (0, -1.52) {
        \textbf{分类 Agent 调度}\\[2pt]
        \texttt{pending} $\rightarrow$ \texttt{classifying} \quad 输出任务代码或空集合
      };

      \node[detect] (detecting) at (0, -3.20) {
        \textbf{子任务检测执行}\\[2pt]
        \texttt{detecting} \quad Reasoning 执行原子检测并上传产物
      };

      \node[summary] (summarizing) at (0, -4.88) {
        \textbf{图像级总结生成}\\[2pt]
        \texttt{summarizing} \quad 基于已落库检测结果生成图像级总结
      };

      \node[state] (converge) at (0, -6.56) {
        \textbf{状态收敛与失败重试}\\[2pt]
        成功进入 \texttt{succeeded} 并开放复核 \quad 失败进入 \texttt{failed} 并等待重试
      };

      \node[review] (review) at (0, -8.24) {
        \textbf{人工复核确认}\\[2pt]
        查看产物、区域指标和 Agent 总结 \quad 选择准确 / 不准确并补充评论
      };

      \node[report] (source) at (0, -9.92) {
        \textbf{项目报告上下文归并}\\[2pt]
        检测事实 + 图像级总结 + 人工复核意见进入报告源数据
      };

      \draw[arrow] (start.south) -- (classifying.north);
      \draw[arrow] (classifying.south) -- (detecting.north);
      \draw[arrow] (detecting.south) -- (summarizing.north);
      \draw[arrow] (summarizing.south) -- (converge.north);
      \draw[arrow] (converge.south) -- (review.north);
      \draw[arrow] (review.south) -- (source.north);

      \node[note, anchor=west] at (5.60, -1.52) {Agent 负责任务选择};
      \node[note, anchor=west] at (5.60, -3.20) {Reasoning 只回调结果};
      \node[note, anchor=west] at (5.60, -6.56) {Core BIZ 统一推进状态};
      \node[note, anchor=west] at (5.60, -8.24) {用户负责复核判断};
      \node[note, anchor=west] at (5.60, -9.92) {报告吸收复核上下文};
    \end{tikzpicture}
  }
  \caption{任务状态流转与人机协同机制流程}\label{fig:design-task-state-human-review-flow}
\end{figure}

状态机与 Agent 方法之间存在直接关系。分类 Agent 的输出决定是否创建子任务，Reasoning 子任务结果决定是否进入图像总结，图像总结结果决定检测主任务是否完成。也就是说，Agent 决策不是孤立文本，而是会改变后续状态流转。因此，状态推进必须由 Core BIZ 统一执行，避免 Activity 服务或前端根据局部信息自行改变业务状态。

在人机协同方面，本文没有把检测结果直接作为最终结论，而是在检测完成后设置人工复核阶段。用户可以同时查看检测产物、区域指标和 Agent 总结，再根据实际情况选择“准确”或“不准确”，也可以补充自己的判断说明。复核信息不会删除原始检测事实，也不会覆盖模型产物，而是作为后续项目报告的一部分输入。这样处理后，系统既保留自动检测形成的证据，也能记录工程人员对结果的判断过程。

人工复核的作用主要体现在两个方面。一方面，幕墙维护建议可能涉及安全责任、维修成本和施工安排，土木工程场景很难完全依赖自动化结论直接决策，因此需要由人对关键结果进行确认。另一方面，复核信息可以为项目级报告提供纠偏依据。若某张图像的检测结果被用户标记为“不准确”，项目报告 Agent 在汇总时应当看到这一信息，并在整体判断中有所体现，而不是把模型输出无条件当作事实。

\section{本章小结}
\label{sec:design-chapter-summary}

本章为建筑幕墙韧性评估软件系统设计了主要技术方案。方案设计中没有把所有智能能力压到一个端到端模型里，而是把任务调度、原子检测、图像级总结、项目报告和人工复核拆成几个边界清楚的环节。分类 Agent 负责把图像语义转换为有限任务代码，Reasoning 服务负责把已有检测模型封装成可调度能力，图像总结 Agent 和项目报告 Agent 再把底层检测结果逐步整理为工程语义。人工复核则保留人的专业判断，使报告结论不只依赖自动检测结果。

这种方案的链路相对较长，但每一步都有对应的证据来源和失败处理位置。相比固定检测流程，它可以根据图像材质选择需要执行的检测能力，减少无意义的模型调用；相比普通工程系统，它又引入 Agent 处理语义判断和信息整合，使系统不再停留于单图单缺陷检测，而是进一步面向项目级韧性评估。\cref{chap:implementation}将在本章基础上继续说明这些设计如何落到多服务架构、协议、数据库、缓存、对象存储、Worker、回调、消息队列、WebSocket、CronJob 和部署配置中。
